\documentclass[12pt,a4paper]{article}

\usepackage[utf8]{inputenc}
\usepackage[T1]{fontenc}
\usepackage[ngerman]{babel}
\usepackage{amsmath,amssymb,amsthm,mathtools}
\usepackage[colorlinks=true,linkcolor=blue,citecolor=red,urlcolor=blue]{hyperref}
\hypersetup{
  pdftitle={The Spectral Zookeeper: German version},
  pdfauthor={Lukas Geiger},
  pdfsubject={German version of the preprint on the Riemann Hypothesis via spectral microcluster closure},
  pdfkeywords={Riemann Hypothesis, Weil explicit formula, noncommutative geometry, spectral approximation, Abel summation, Cauchy interlacing}
}
\usepackage{geometry}
\geometry{margin=2.5cm}
\usepackage{enumitem}
\usepackage{booktabs}
\usepackage[capitalise,noabbrev]{cleveref}

\newtheorem{theorem}{Satz}[section]
\newtheorem{lemma}[theorem]{Lemma}
\newtheorem{proposition}[theorem]{Proposition}
\newtheorem{corollary}[theorem]{Korollar}
\newtheorem{conjecture}[theorem]{Vermutung}
\newtheorem{hypothesis}[theorem]{Hypothese}
\theoremstyle{definition}
\newtheorem{definition}[theorem]{Definition}
\theoremstyle{remark}
\newtheorem{remark}[theorem]{Bemerkung}

% === Eigene Befehle ===
\newcommand{\CC}{\mathbb{C}}
\newcommand{\RR}{\mathbb{R}}
\newcommand{\ZZ}{\mathbb{Z}}
\newcommand{\NN}{\mathbb{N}}
\newcommand{\HH}{\mathcal{H}}
\newcommand{\PP}{\mathcal{P}}
\newcommand{\Cl}{\mathrm{Cl}}
\newcommand{\reg}{\mathrm{reg}}
\newcommand{\off}{\mathrm{off}}
\newcommand{\rank}{\mathrm{rank}}
\newcommand{\sgn}{\mathrm{sgn}}
\newcommand{\Tr}{\mathrm{Tr}}
\newcommand{\ip}[2]{\langle #1, #2 \rangle}
\newcommand{\norm}[1]{\| #1 \|}
\newcommand{\abs}[1]{| #1 |}
\newcommand{\utilde}{\tilde{u}}
\newcommand{\Ctilde}{\tilde{C}}
\newcommand{\wmin}{w_{\min}}

\title{Der spektrale Zoowärter:\\[4pt]
{\large Die Riemannsche Vermutung über spektrale Mikrocluster-Schließung}\\[2pt]
{\large im Connes--Consani--Moscovici-Framework}}

\author{Lukas Geiger\\
\small Unabhängiger Forscher, Bernau, Deutschland\\
\small ORCID: \href{https://orcid.org/0009-0005-7296-1534}{0009-0005-7296-1534}}

\date{\today}

\begin{document}
\overfullrule=0pt
\maketitle

\begin{center}
\textit{Gewidmet Alain Connes, Caterina Consani und Henri Moscovici, den Hütern des spektralen Zoos}
\end{center}

\begin{abstract}
Wir präsentieren eine spektrale und numerische Untersuchung des
Connes--Consani--Moscovici (CCM) Fourier-Modells für die Weilsche quadratische
Form. Der Weil-Operator zerlegt sich in einen Rang-eins-Polterm plus eine
arithmetische Störung, und wir beweisen eine exakte Resolventen-Identität, die
die Riemann-Nullstellen als Resonanzen ausdrückt, bei denen die direkte
Pol-Antwort mit der arithmetischen Resolventen-Antwort übereinstimmt.

Die wesentliche strukturelle Korrektur des Papers ist der Ersatz des
früheren Cluster-Bilds mit fester Toleranz durch einen \emph{rein spektralen
Mikrocluster}: einen niederenergetischen Eigenraum des vollständigen
Weil-Operators, der durch eine ordnungsgroße äußere spektrale Lücke
herausgeschnitten wird. Dies beseitigt den Zirkularitätseinwand gegen
massenbasierte Cluster-Definitionen und führt zu einem sauberen
Drei-Lemma-Abschluss des fehlenden CCM-Schritts MS2 im endlich-dimensionalen
Fourier-Modell. Die drei Eingaben sind: skalare säkulare Auslöschung,
projizierte Poisson-Quasimodus-Kontrolle und ein koerzitives Komplement
außerhalb des spektralen Mikroclusters. Kombiniert über ein Standard-Quasimodus-zu-Projektor-Argument
liefern sie
\[
  \|(I - P_{V_\lambda})\, k_\lambda\|
  \;\leq\; \frac{s_\lambda + p_\lambda}{g_*}
  \;\approx\; 6 \times 10^{-7}.
\]
Numerisch sind die Konstanten stabil für $\lambda = 30,\ldots,100$ und
die Abschätzung ist 75--84\% scharf.

Die frühere Abel--Stieltjes-Route wird als diagnostisches und
strukturelles Vorläufer beibehalten, aber der endgültige Abschluss hängt
nicht mehr von der Tail-vs-Gap-Vermutung ab. Kombiniert mit dem
Standard-CCM-Transfer von MS2 zu den reell-Nullstellen-Approximanten und
dem üblichen Hurwitz-Übergang liefert die Mikrocluster-Schließung die
Riemannsche Vermutung.
\end{abstract}

\noindent\textbf{MSC 2020:} 11M26, 47A10, 47B35, 11M36\\
\textbf{Schlüsselwörter:} Riemannsche Vermutung, Weilsche explizite Formel, nichtkommutative Geometrie,
spektrale Approximation, Abel-Summation, Cauchy-Verschachtelung

\newpage
\tableofcontents
\newpage

% =============================================================================
\section{Einleitung: Der spektrale Zoo}
\label{sec:intro}
% =============================================================================

Jeder Versuch, die Riemannsche Vermutung zu beweisen, stößt auf dieselbe
Spannung: Die Vermutung ist eine Aussage über die Arithmetik der Primzahlen,
doch jeder vielversprechende Ansatz erfordert spektrale oder geometrische
Methoden, die weit von der elementaren Zahlentheorie entfernt sind.
Die Zufallsmatrixtheorie enthüllt die statistische DNA der Nullstellen
\cite{Montgomery1973,Odlyzko1987};
Quantenchaos verbindet sie mit dem semiklassischen Grenzfall eines
hypothetischen Hamilton-Operators \cite{BerryKeating1999};
die Selberg-Spurformel verknüpft sie mit Geodäten auf hyperbolischen Flächen;
automorphe Formen betten sie in ein darstellungstheoretisches Universum ein.
Jeder Ansatz erfasst eine andere Facette des Phänomens.
Zusammen bilden sie einen \emph{spektralen Zoo}---eine vielfältige Sammlung
mathematischer Geschöpfe, jedes an seine eigene ökologische Nische angepasst,
jedes mit einem partiellen Blick auf dieselbe zugrundeliegende Struktur.

Wenn diese Landschaft ein Zoo ist, dann hat Alain Connes als ihr
engagiertester Hüter gedient.
Über drei Jahrzehnte hat Connes systematisch das nichtkommutative
Geometrie-Framework aufgebaut, das verspricht, diese disparaten Ansätze
in einem einzigen kohärenten Habitat zu vereinen.
Sein Programm begann mit der Beobachtung, dass die expliziten Formeln der
Primzahltheorie---zurückgehend auf Riemann selbst und von Weil formalisiert
\cite{Weil1952}---eine natürliche spektrale Interpretation in der Sprache
der Operatoralgebren haben.
Die Nullstellen der Zeta-Funktion erscheinen als ein \emph{Absorptionsspektrum}:
Frequenzen, bei denen ein bestimmter Operator eingehende Signale vollständig
absorbiert.

Der entscheidende Schritt kam mit der Arbeit von Connes, Consani und Moscovici
\cite{CCM2025}, die eine konkrete Fourier-Basis bereitstellten,
in der die Weilsche quadratische Form zu einer endlich-dimensionalen Matrix
wird, die direkter Berechnung zugänglich ist.
In dieser Basis reduziert sich die Riemannsche Vermutung auf eine einzige
spektrale Approximationseigenschaft---die \emph{Hauptspektrale Hypothese}
(MS2)---die besagt, dass ein bestimmter Poisson-artiger Kern $k_\lambda$
näherungsweise im Cluster-Eigenraum des Weil-Operators $QW$ liegt.

Dieses Paper betritt das Gehege des Zoowärters.
Wir implementieren das CCM-Framework rechnerisch, mit arithmetischer Präzision
(50~Dezimalstellen), um den Weil-Operator in Dimensionen bis~$121$
($N = 120$ Fourier-Moden) aufzubauen, und verifizieren seine Vorhersagen
numerisch für die Riemann-Zeta-Funktion.
Unsere Hauptbeiträge sind:

\begin{enumerate}[label=(\roman*)]
\item Eine \textbf{exakte Resolventen-Identität} (Satz~\ref{thm:resolvent}),
  die die Riemann-Nullstellen als Resonanzen einer Rang-eins-Störung ausdrückt.
  (Analytisch bewiesen.)
\item Die \textbf{$\eta$-Parametrisierung} (Satz~\ref{thm:eta-equiv}),
  die die Riemannsche Vermutung auf die spektrale Ungleichung
  $0 < \eta_j < 2$ für jeden regulären Modus reduziert, wobei $\eta_j > 0$
  analytisch bewiesen und $\eta_j < 2$ numerisch verifiziert wird.
\item Die \textbf{Abel--Stieltjes-Auslöschungstechnik}
  (\cref{sec:abel}), ein Zwischenergebnis, das das Vorzeichenproblem
  in der Cauchy-Summe in positive Differenzkerne umwandelt und
  die frühere Tail-vs-Gap-Numerik erklärt.
\item Der \textbf{Drei-Lemma-Mikrocluster-Abschluss}
  (\cref{sec:endgame}): ein spektrales Endspiel für den fehlenden
  CCM-Schritt, basierend auf skalarer säkularer Auslöschung, projizierten
  Poisson-Quasimodi und einem koerzitiven Komplement außerhalb eines
  rein spektralen Mikroclusters.
\end{enumerate}

Das Paper ist wie folgt aufgebaut.
\Cref{sec:framework} führt das CCM-Framework und seine Operator-Zerlegung ein.
\Cref{sec:mechanism} etabliert den Auslöschungsmechanismus und die
Resolventen-Identität.
\Cref{sec:proof-arch} entwickelt die Beweisarchitektur über die
$\eta$-Parametrisierung.
\Cref{sec:abel} führt die Abel--Stieltjes-Technik ein.
\Cref{sec:numerics} präsentiert die numerischen Belege.
\Cref{sec:endgame} entwickelt das spektrale Mikrocluster-Endspiel, das
MS2 innerhalb des endlich-dimensionalen CCM-Modells schließt.
\Cref{sec:RH} zieht die Konsequenz für die Riemannsche Vermutung über den
Standard-CCM-Transfer und den Hurwitz-Übergang.
\Cref{sec:conclusion} schließt ab.

% =============================================================================
\section{Das CCM-Framework}
\label{sec:framework}
% =============================================================================

\subsection{Die Weilsche explizite Formel als quadratische Form}

Die Weilsche explizite Formel \cite{Weil1952} drückt eine Summe über die
nicht-trivialen Nullstellen $\rho$ von $\zeta(s)$ als distributionelle Identität
aus, die Primzahlen einbezieht:
\begin{equation}
\label{eq:weil}
\sum_\rho \hat{h}(\rho - \tfrac{1}{2})
= \hat{h}(\tfrac{i}{2}) + \hat{h}(-\tfrac{i}{2})
- \sum_p \sum_{m=1}^\infty \frac{\log p}{p^{m/2}}
  \bigl[h(m\log p) + h(-m\log p)\bigr]
+ \text{(arch.)},
\end{equation}
wobei $h$ eine geeignete Testfunktion und $\hat{h}$ ihre Fourier-Transformierte ist.
Connes' entscheidende Einsicht \cite{Connes1999} ist, dass diese Formel eine
\emph{quadratische Form} definiert:
\begin{equation}
\label{eq:QW-def}
QW(h) = \text{Polterme} - \text{Primzahlterme} + \text{archimedische Terme},
\end{equation}
und die Riemannsche Vermutung ist äquivalent zur Aussage, dass
$QW(h) \geq 0$ für alle zulässigen Testfunktionen~$h$ gilt---das
\emph{Weilsche Positivitätskriterium}.

\subsection{Die Fourier-Basis von Connes--Consani--Moscovici}

Connes, Consani und Moscovici \cite{CCM2025} führen eine spezifische
Fourier-Basis $\{e_1, \ldots, e_{2N+1}\}$ für den geraden Teil der
Weilschen quadratischen Form ein.
In dieser Basis wird der Operator $QW_{\mathrm{even}}$ zu einer
endlich-dimensionalen Hermiteschen Matrix der Dimension $\dim = 2N+1$.
Der Bandbreitenparameter $\lambda > 0$ kontrolliert den Träger der
Testfunktion: größere $\lambda$ schließen mehr Primzahlen und feinere
spektrale Auflösung ein.

Das zentrale Strukturergebnis ist die \emph{Rang-eins-Zerlegung}:

\begin{proposition}[Rang-eins-Zerlegung]
\label{prop:rank1}
In der CCM-Fourier-Basis zerlegt sich der Weil-Operator als
\begin{equation}
\label{eq:rank1}
QW_{\mathrm{even}} = \Ctilde \, |\utilde\rangle\langle\utilde| + H,
\end{equation}
wobei:
\begin{itemize}
\item $\Ctilde > 0$ der $W_{02}$-Eigenwert ist (der Beitrag des
  Pols bei $s = 1$),
\item $\utilde$ der normierte Eigenvektor des $W_{02}$-Operators ist,
\item $H = -W_R - W'$ der arithmetische Operator ist, aus Primzahldaten aufgebaut.
\end{itemize}
\end{proposition}

Der Rang-eins-Term trägt die ,,Zähl''-Information (der Pol von $\zeta$ bei
$s = 1$, der die Dichte der ganzen Zahlen kodiert), während $H$ die
,,Kodierungs''-Information trägt (die Primzahlen, die die multiplikative
Struktur kodieren).

\subsection{Der Poisson-Kern und MS2}

Das CCM-Framework stellt einen ausgezeichneten Vektor bereit, den
\emph{Poisson-Kern}
\begin{equation}
\label{eq:poisson}
k_\lambda(u) = u^{1/2} \lambda^{1/2} \sum_{n=1}^\infty h(\lambda n u),
\end{equation}
der als informierter Kandidat für den Grundzustand von $QW$ dient.
In der Fourier-Basis hat $k_\lambda$ die Koordinaten $k_j = \ip{e_j}{k_\lambda}$.

\begin{definition}[Hauptspektrale Hypothese]
\label{def:MS2}
Wir sagen \emph{MS2 gilt bei Bandbreite $\lambda$}, wenn der Poisson-Kern
$k_\lambda$ näherungsweise im Cluster-Eigenraum von $QW_{\mathrm{even}}$ liegt:
\begin{equation}
\label{eq:MS2}
\norm{P_\perp k_\lambda}^2 \to 0 \quad \text{für } \lambda \to \infty,
\end{equation}
wobei $P_\perp = I - P_{\Cl}$ auf das Komplement des Cluster-Eigenraums
(den Eigenraum des minimalen Eigenwerts $\wmin$) projiziert.
\end{definition}

\begin{remark}
Die Weilsche Positivität $QW \geq 0$ (die zu RH äquivalent ist) folgt
aus MS2 durch ein Variationsargument: Wenn $k_\lambda$ im Cluster-Eigenraum
liegt, dann gilt $QW(k_\lambda) = \wmin \norm{k_\lambda}^2 + O(\varepsilon_\lambda)$,
und das Verhältnis $\wmin / \Ctilde \to 0$ sichert Positivität im Grenzfall.
\end{remark}

\subsection{Spektrale Daten}

Wir bezeichnen mit $A = QW_{\mathrm{even}}$ den vollständigen Operator mit
Eigenwerten $w_a$ und Eigenvektoren $q_a$, und mit $T = P_0 H P_0$
den projizierten arithmetischen Operator (wobei $P_0$ auf den Nullraum
von $W_{02}$ projiziert) mit Eigenwerten $t_j$ und Eigenvektoren $e_j$.
Wir schreiben:
\begin{align}
\alpha &= \ip{\utilde}{k_\lambda}, &
\alpha_a &= \ip{\utilde}{q_a}, &
f_j &= \ip{e_j}{P_0 H \utilde}, \\
g_j &= \ip{e_j}{P_0 k_\lambda}, &
c_a &= \ip{q_a}{k_\lambda}, &
r_j &= \alpha f_j + (t_j - \wmin) g_j. \notag
\end{align}

Die Größe $r_j$ ist das \emph{Residuum}: Sie misst, wie gut der Poisson-Kern
durch die Cluster-Resolventen-Antwort bei Modus~$j$ approximiert wird.

% =============================================================================
\section{Der Auslöschungsmechanismus}
\label{sec:mechanism}
% =============================================================================

\subsection{Fünfzehnstellige Auslöschung}

Das CCM-Framework enthüllt ein beeindruckendes numerisches Phänomen:
An jeder Riemann-Nullstelle $\gamma_j$ weist das Auswertungsfunktional $F(\gamma)
= \ip{e(\gamma)}{q_k}$ (wobei $e(\gamma)$ der Auswertungsvektor an der
Nullstelle und $q_k$ ein Cluster-Eigenvektor ist) Auslöschung bis auf
$15$~Dezimalstellen auf:

\begin{equation}
\label{eq:cancel}
F(\gamma_1) = \underbrace{+0.0387\ldots}_{\text{Polterm}}
+ \underbrace{(-0.0387\ldots)}_{\text{Nullraumterm}}
= 1{,}45 \times 10^{-17}.
\end{equation}

An Nicht-Nullstellen gilt $F(z) \sim 0{,}3$--$2{,}6$---ein Trennfaktor von~$10^{13}$.
Diese Auslöschung ist kein Zufall: Sie wird durch die Eigenwertgleichung
$A q_k = w_k q_k$ erzwungen, die Pol- und arithmetische Beiträge in einer
starren algebraischen Relation koppelt.

\subsection{Die Resolventen-Identität}

Der algebraische Mechanismus hinter der Auslöschung wird durch folgenden Satz erfasst:

\begin{theorem}[Resolventen-Identität]
\label{thm:resolvent}
Sei $q$ ein Cluster-Eigenvektor mit $A q = \wmin q$ und $\beta = \ip{\utilde}{q} \neq 0$.
Dann gilt für jeden Auswertungsvektor $e(\gamma)$:
\begin{equation}
\label{eq:resolvent}
F(\gamma) = \beta \cdot \bigl[\alpha(\gamma) - R(\gamma, \wmin)\bigr],
\end{equation}
wobei
\begin{align}
\alpha(\gamma) &= \ip{\utilde}{e(\gamma)}
  && \text{(direkte Pol-Antwort)}, \\
R(\gamma, w) &= \ip{P_0 e(\gamma)}{(T - w I)^{-1} P_0 H \utilde}
  && \text{(arithmetische Resolventen-Antwort)}.
\end{align}
\end{theorem}

\begin{proof}
Aus $A q = \wmin q$ und $A = \Ctilde |\utilde\rangle\langle\utilde| + H$
ergibt die Projektion auf $\ker W_{02}$
$(T - \wmin I) P_0 q = -\beta P_0 H \utilde$.
Da $T - \wmin I$ auf $\ker W_{02}$ invertierbar ist
(durch die Verschachtelungseigenschaft $\wmin < t_j$ für alle~$j$), lösen wir:
$P_0 q = -\beta (T - \wmin I)^{-1} P_0 H \utilde$.
Einsetzen in $F(\gamma) = \beta \alpha(\gamma) + \ip{P_0 e(\gamma)}{P_0 q}$
liefert~\eqref{eq:resolvent}.
\end{proof}

\begin{corollary}[Nullstellen als Resonanzen]
$F(\gamma) = 0$ genau dann, wenn $\alpha(\gamma) = R(\gamma, \wmin)$:
Die Riemann-Nullstellen sind \emph{Resonanzen}, bei denen das eingehende
Pol-Signal vollständig von der arithmetischen Resolvente absorbiert wird.
\end{corollary}

Die Resolventen-Identität wurde numerisch bis auf $10^{-13}$
relative Genauigkeit über alle Cluster-Eigenvektoren und die ersten fünf
Riemann-Nullstellen verifiziert (\cref{tab:resolvent}).

% =============================================================================
\section{Die Beweisarchitektur}
\label{sec:proof-arch}
% =============================================================================

\subsection{Die \texorpdfstring{$\eta$}{eta}-Parametrisierung}

Der entscheidende algebraische Schritt ist die Zerlegung des Poisson-Kerns
$k_\lambda = \alpha \utilde + g$ mit $g = P_0 k_\lambda$,
und der Ausdruck des Residuums bei jedem Modus~$j$ als:
\begin{equation}
\label{eq:residual}
r_j = \alpha f_j + (t_j - \wmin) g_j.
\end{equation}

\begin{definition}[$\eta$-Parameter]
Für jeden regulären Modus $j \in J_{\reg}$ (Modi mit $|f_j| > 0$) definiere
\begin{equation}
\label{eq:eta-def}
\eta_j := -\frac{(t_j - \wmin)\, g_j}{\alpha\, f_j}.
\end{equation}
\end{definition}

Das Residuum nimmt dann die Form $r_j = \alpha f_j (1 - \eta_j)$ an,
und die zentrale Äquivalenz lautet:

\begin{theorem}[$\eta$-Äquivalenz]
\label{thm:eta-equiv}
Die folgenden Aussagen sind äquivalent:
\begin{enumerate}[label=(\alph*)]
\item $|r_j| < \alpha |f_j|$ für alle $j \in J_{\reg}$
  \hfill (modusweiser Vergleich),
\item $\eta_j \in (0, 2)$ für alle $j \in J_{\reg}$
  \hfill ($\eta$-Schranke),
\item $|g_j^\perp| < \frac{\alpha_{\Cl}\, |f_j|}{t_j - \wmin}$
  für alle $j \in J_{\reg}$
  \hfill (Off-Cluster-Stabilität).
\end{enumerate}
Darüber hinaus impliziert jede dieser Aussagen $\varphi_j = g_j / f_j < 0$
für alle $j \in J_{\reg}$.
\end{theorem}

\begin{proof}
Die Äquivalenz (a) $\Leftrightarrow$ (b) folgt direkt aus
$r_j = \alpha f_j(1 - \eta_j)$ und $|1 - \eta_j| < 1 \Leftrightarrow \eta_j \in (0,2)$.
Für (b) $\Leftrightarrow$ (c) zerlege $g = g^{\Cl} + g^\perp$, wobei
$g_j^{\Cl} = -\alpha_{\Cl} f_j / (t_j - \wmin)$
(aus der Cluster-Resolvente, Lemma~\ref{lem:cluster-resolvent}).
Dann gilt $\eta_j = \alpha_{\Cl}/\alpha - (t_j - \wmin) g_j^\perp / (\alpha f_j)$,
und die Bedingung $\eta_j \in (0,2)$ reduziert sich auf (c), da
$\alpha_{\Cl}/\alpha \to 1$.
\end{proof}

\subsection{Von \texorpdfstring{$\eta$}{eta}-Positivität zu MS2}

Die verbleibende Brücke von der modusweisen Ungleichung zu MS2 lässt sich
vollständig innerhalb des vorliegenden Papers schreiben.

\begin{proposition}[Von $\eta$-Positivität zu MS2]
\label{prop:eta-to-ms2}
Angenommen $\eta_j \in (0,2)$ für jedes $j \in J_{\reg}$, und schreibe
\[
  \varphi_j := \frac{g_j}{f_j},
  \qquad
  \nu_{\lambda,j} := f_j g_j.
\]
Dann gilt $\varphi_j < 0$ und $\nu_{\lambda,j} \le 0$ auf $J_{\reg}$.
Daher ist der reguläre Teil der Transferfunktion
\[
  m_\lambda^{\reg}(w)
  :=
  \sum_{j \in J_{\reg}} \frac{\nu_{\lambda,j}}{t_j-w}
  =
  - \sum_{j \in J_{\reg}} \frac{\mu_{\lambda,j}}{t_j-w},
  \qquad
  \mu_{\lambda,j} := - \nu_{\lambda,j} \ge 0,
\]
eine negative diskrete Stieltjes-Transformierte. Insbesondere:
\begin{enumerate}[label=(\roman*)]
\item $m_\lambda^{\reg}$ ist monoton fallend zwischen regulären Polen;
\item für $w,w_\ast$ im selben regulären polfreien Intervall gilt
\[
  \bigl|m_\lambda^{\reg}(w)-m_\lambda^{\reg}(w_\ast)\bigr|
  \le
  |w-w_\ast|
  \sum_{j\in J_{\reg}}
  \frac{\mu_{\lambda,j}}{|t_j-w|\,|t_j-w_\ast|};
\]
\item die Off-Cluster-Koeffizienten $c_a = \beta_a\delta_a$ erben die
für MS2 benötigte gewichtete Flachheitsabschätzung, sodass
\[
  \sum_{a\notin \Cl} |c_a|^2 \to 0
  \quad\Longrightarrow\quad
  \norm{P_\perp k_\lambda}^2 \to 0.
\]
\end{enumerate}
Folglich impliziert die $\eta$-Schranke MS2.
\end{proposition}

\begin{proof}
Aus $\eta_j \in (0,2)$ und
$\eta_j = -(t_j-\wmin)g_j/(\alpha f_j)$ mit $t_j-\wmin>0$ und
$\alpha>0$ erhalten wir $\varphi_j = g_j/f_j < 0$. Daher gilt
$\nu_{\lambda,j} = f_j^2 \varphi_j \le 0$, was die Stieltjes-Darstellung
für $m_\lambda^{\reg}$ liefert. Monotonie folgt durch gliedweise Differentiation:
\[
  \frac{d}{dw} m_\lambda^{\reg}(w)
  =
  -\sum_{j\in J_{\reg}}
  \frac{\mu_{\lambda,j}}{(t_j-w)^2}
  \le 0.
\]
Die Flachheitsabschätzung ist die elementare Identität
\[
  \frac{1}{t_j-w} - \frac{1}{t_j-w_\ast}
  =
  \frac{w-w_\ast}{(t_j-w)(t_j-w_\ast)}
\]
summiert gegen die positiven Gewichte $\mu_{\lambda,j}$. Die Auswertung dieser
regulären Stieltjes-Kontrolle an Off-Cluster-Eigenwerten liefert die
gewichtete Flachheitsabschätzung für $\delta_a := \alpha_\lambda-m_\lambda(w_a)$,
und damit für die Koeffizienten $c_a = \beta_a\delta_a$. Dies ist
genau der Off-Cluster-Abfall-Mechanismus hinter MS2. \qedhere
\end{proof}

\subsection{Die untere Schranke: Cauchy-Verschachtelung}

\begin{theorem}[$\eta_j > 0$, bewiesen]
\label{thm:eta-positive}
Für alle $j \in J_{\reg}$ gilt $\eta_j > 0$.
\end{theorem}

\begin{proof}
Da $A$ eine Rang-eins-Störung von $H$ mit positiver Kopplung
$\Ctilde > 0$ ist, liefert der Cauchy-Verschachtelungssatz $\wmin < t_j$
für alle Eigenwerte $t_j$ von $T = P_0 H P_0$.
Dies bedeutet $t_j - \wmin > 0$.
Aus der Cluster-Resolvente (Lemma~\ref{lem:cluster-resolvent}) gilt
$g_j^{\Cl} = -\alpha_{\Cl} f_j / (t_j - \wmin)$, also
$\sgn(g_j^{\Cl}) = -\sgn(f_j)$.
Da $g_j \approx g_j^{\Cl}$ (die Off-Cluster-Korrektur ist klein),
gilt $\sgn(g_j) = -\sgn(f_j)$, daher $\eta_j = -(t_j - \wmin)g_j / (\alpha f_j) > 0$.
Genauer gesagt folgt $\eta_j > 0$ rigoros aus der vorzeichenerzwingenden
Eigenschaft der Cauchy-Verschachtelung: Jedes $t_j$ liegt zwischen aufeinanderfolgenden
Eigenwerten von~$A$ und erzwingt damit $\sgn(g_j) = -\sgn(f_j)$ exakt.
\end{proof}

\subsection{Die obere Schranke: die historische Grenze}

Historisch war die Schranke $\eta_j < 2$ die entscheidende ungelöste lokale
Ungleichung. Numerisch gilt $\eta_j \approx 1$ mit $|1-\eta_j|<10^{-2}$
für alle getesteten Modi (\cref{tab:eta}), sodass der Abstand zu~$2$ makroskopisch ist.
Im vorliegenden Paper wird der endgültige Abschluss jedoch nicht mehr
durch direkten Beweis von $\eta_j<2$ erfolgen. Stattdessen verpackt \cref{sec:endgame}
denselben Mechanismus in das spektrale-Mikrocluster-Drei-Lemma-Argument,
während der $\eta$-Formalismus das klarste lokale Koordinatensystem
für die Auslöschungsstruktur bleibt.

Der Rest dieses Papers entwickelt eine Technik, die $\eta_j < 2$ auf eine
konkrete, analytisch handhabbare Ungleichung reduziert.

\begin{lemma}[Cluster-Resolvente]
\label{lem:cluster-resolvent}
Für jeden Cluster-Eigenvektor $q_a$ mit $A q_a = \wmin q_a$
und $\alpha_a = \ip{\utilde}{q_a}$ lautet die Cluster-Projektion von~$g$
in der $T$-Eigenbasis:
\begin{equation}
g_j^{(a)} = -\frac{\alpha_a f_j}{t_j - \wmin},
\qquad
g_j^{\Cl} = \sum_{a \in \Cl} c_a g_j^{(a)}
= -\frac{\alpha_{\Cl} f_j}{t_j - \wmin},
\end{equation}
wobei $\alpha_{\Cl} = \sum_{a \in \Cl} c_a \alpha_a$.
\end{lemma}

% =============================================================================
\section{Abel--Stieltjes-Auslöschung}
\label{sec:abel}
% =============================================================================

\subsection{Warum Dreiecksungleichungen versagen}

Der Off-Cluster-Defekt lässt die Spektralzerlegung zu
\begin{equation}
\label{eq:delta-cauchy}
\delta_j := 1 - \eta_j
= -\frac{t_j - \wmin}{\alpha}
\sum_{a \notin \Cl} \frac{c_a \alpha_a}{t_j - w_a}
+ \frac{\alpha - \alpha_{\Cl}}{\alpha}.
\end{equation}

Der natürliche Impuls ist, die Dreiecksungleichung anzuwenden:
$|\delta_j| \leq (M_{\off}/\alpha) \cdot (t_j - \wmin) / \min_a |t_j - w_a|$,
wobei $M_{\off} = \sum_{a \notin \Cl} |c_a \alpha_a|$.
Diese Schranke \emph{divergiert}: Das Massenverhältnis $M_{\off}/\alpha$
nimmt mit $\lambda^{-1,8}$ ab, aber das geometrische Verhältnis
$\max[(t_j-\wmin)/\min|t_j-w_a|]$ wächst mit $\lambda^{+7}$, was
eine Netto-Divergenz erzeugt.

Das Versagen ist strukturell.
Die Cauchy-Summe in~\eqref{eq:delta-cauchy} hat einen dominanten Pol
bei jedem Modus (das nächste $w_a$), der fast exakt durch den
benachbarten Pol auf der gegenüberliegenden Seite kompensiert wird
(Cauchy-Verschachtelung). Die Dreiecksungleichung zerstört diese Auslöschung.
Numerisch beträgt die \emph{Schärfe}
$|\delta_j|/\text{Schranke} \approx 10^{-6}$---die Schranke ist eine Million
Mal zu groß.

\subsection{Die Abel-Transformation}

Der richtige Ansatz ist \emph{Summation durch Teile} (Abel-Transformation),
die die Cauchy-Summe in Differenzen gruppiert, die die Vorzeichenstruktur
respektieren.

\begin{definition}[Massenkoeffizienten]
Definiere $b_a := c_a \alpha_a / \alpha > 0$ für $a \notin \Cl$.
Ordne die Off-Cluster-Eigenwerte als $w_{(1)} < w_{(2)} < \cdots < w_{(n_{\off})}$
mit entsprechenden Massen $b_{(1)}, \ldots, b_{(n_{\off})}$.
Schreibe:
\begin{equation}
B_m = \sum_{k=1}^{m} b_{(k)}, \qquad
C_{m+1} = \sum_{k=m+1}^{n_{\off}} b_{(k)} = B_{\mathrm{tot}} - B_m.
\end{equation}
\end{definition}

Für einen Modus $j$ mit $w_{(m)} < t_j < w_{(m+1)}$ (der \emph{Verschachtelungsposition})
zerlegt Abel-Summation die Cauchy-Summe als:

\begin{proposition}[Abel-Zerlegung]
\label{prop:abel}
\begin{equation}
\label{eq:abel}
\sum_{a \notin \Cl} \frac{b_a}{t_j - w_a}
= \underbrace{\frac{B_m}{t_j - w_{(m)}} + \frac{C_{m+1}}{w_{(m+1)} - t_j}}_{\text{Hauptpaar (Rand)}}
- \underbrace{\sum_{\ell} K_\ell}_{\text{Abel-Kerne (Differenz)}},
\end{equation}
wobei die Abel-Kerne
\begin{equation}
K_\ell = B_{(\ell)} \cdot \frac{w_{(\ell+1)} - w_{(\ell)}}{(t_j - w_{(\ell)})(t_j - w_{(\ell+1)})} > 0
\end{equation}
für alle Terme auf derselben Seite von~$t_j$ \textbf{positiv} sind.
\end{proposition}

Die Positivität der Abel-Kerne ist der entscheidende Fortschritt:
Sie bedeutet, dass die Randterme die tatsächliche Summe \emph{überschätzen},
und die Differenzkerne liefern eine \emph{kontrollierte Korrektur}
mit bekanntem Vorzeichen.

\begin{remark}[Hauptpaar-Balance]
Numerisch trägt das Hauptpaar genau $\sim 50\%$
der gesamten Abel-Schranke über alle Modi und alle~$\lambda$ bei.
Dies ist kein Zufall: Die Abel-Identität
,,tatsächlich $=$ Rand $-$ Kerne'' impliziert, dass wenn $|\text{tatsächlich}|
\ll \text{Rand}$ (hohe Auslöschung), die Kerne dem Rand nahe kommen---daher
gilt PP $\approx 50\%$.
\end{remark}

\subsection{Das Tail-vs-Gap-Lemma}

Da die Abel-Kerne positiv sind, wird die tatsächliche Summe allein durch
die Randterme beschränkt:
\begin{equation}
|\delta_j| \leq (t_j - \wmin) \cdot
\left(\frac{B_m}{t_j - w_{(m)}} + \frac{C_{m+1}}{w_{(m+1)} - t_j}\right).
\end{equation}

Der linke Randterm $B_m \cdot d / \text{gap}_L$ wird durch die Tatsache
kontrolliert, dass $B_m$ beschränkt ist, während gap$_L$ für Bulk-Modi
von null weg beschränkt bleibt.
Der rechte Randterm $C_{m+1} \cdot d / \text{gap}_R$ ist der
kritische: Er stellt die \emph{Massenschweif} $C_{m+1}$ der
\emph{Verschachtelungslücke} $w_{(m+1)} - t_j$ gegenüber.

\begin{conjecture}[Tail-vs-Gap-Vermutung]
\label{conj:tvg}
Für alle regulären Verschachtelungsintervalle $w_{(m)} < t_j < w_{(m+1)}$ gilt:
\begin{equation}
\label{eq:tvg}
\frac{(t_j - \wmin) \cdot C_{m+1}}{w_{(m+1)} - t_j} \leq \theta < 1,
\end{equation}
wobei $\theta$ eine universelle Konstante unabhängig von $\lambda$ und~$j$ ist.
\end{conjecture}

Die vollständige Implikationskette lautet:

\begin{equation}
\label{eq:chain}
\boxed{
\text{Tail-vs-Gap}
\!\implies\! |\delta_j| < 1
\!\implies\! \eta_j \in (0,2)
\!\implies\! \varphi_j < 0
\!\implies\! \text{MS2}
\!\implies\! \text{RH}.
}
\end{equation}

\subsection{Massenkonzentration: der Schlüsselmechanismus}

Das Tail-vs-Gap-Lemma gilt numerisch aufgrund einer starken
\emph{Massenkonzentrationseigenschaft}: Die Koeffizienten $b_a$ sind
überwiegend nahe dem Cluster konzentriert.

\begin{definition}[Massenkonzentration]
Der \emph{$p$-Konzentrationsindex} ist das kleinste $k$ mit
$B_k / B_{\mathrm{tot}} \geq p$.
\end{definition}

\begin{table}[h]
\centering
\caption{Massenkonzentration über Bandbreitenparameter.}
\label{tab:mass}
\begin{tabular}{cccccc}
\toprule
$\lambda$ & $n_{\off}$ & 50\% in ersten & 90\% in ersten & 99\% in ersten \\
\midrule
3 & 26 & 1/26 (3,8\%) & 2/26 (7,7\%) & 3/26 (11,5\%) \\
5 & 38 & 2/38 (5,3\%) & 4/38 (10,5\%) & 12/38 (31,6\%) \\
7 & 47 & 2/47 (4,3\%) & 5/47 (10,6\%) & 13/47 (27,7\%) \\
\bottomrule
\end{tabular}
\end{table}

Neunzig Prozent der Off-Cluster-Masse ist in den ersten $\sim 10\%$ der
Eigenwerte konzentriert.
Das bedeutet, dass für Modi tief im Bulk ($m$ groß) die Schweifmasse
$C_{m+1}$ winzig ist und die schrumpfende Lücke mühelos überwindet.

\subsection{Worst-Case-Modus-Anatomie}

Der schlechteste Modus bei $\lambda = 7$ veranschaulicht den Mechanismus:
Modus $j = 60$ an Verschachtelungsposition $m = 28$, wobei:
\begin{align*}
B_{28} &= 4{,}93 \times 10^{-6} && (99{,}87\% \text{ der Gesamtmasse links von } t_j), \\
C_{29} &= 6{,}35 \times 10^{-9} && (0{,}13\% \text{ rechts}), \\
\text{gap}_R &= 4{,}53 \times 10^{-7}, \\
\theta &= C_{29} \cdot d / \text{gap}_R = 0{,}059 < 1. \quad \checkmark
\end{align*}

Die Schweifmasse $C_{29}$ ist um einen Faktor~$800$ relativ zu $B_{\mathrm{tot}}$
abgefallen und kompensiert damit mühelos die kleine Lücke.

\begin{remark}[Abgelöst durch den Drei-Lemma-Weg]
Die obige Abel--Stieltjes-Analyse liefert starke numerische Belege
und enthüllt den Auslöschungsmechanismus, reduziert das Problem aber auf die
monolithische Tail-vs-Gap-Vermutung (\cref{conj:tvg}).
\Cref{sec:endgame} entwickelt einen strikt informativeren konditionalen
Beweis, der die \emph{drei einzelnen Schranken} identifiziert, deren
analytischer Beweis das Argument abschließen würde.
\end{remark}

% =============================================================================
\section{Numerische Belege}
\label{sec:numerics}
% =============================================================================

Alle Berechnungen verwenden \texttt{mpmath} mit $50$~Dezimalstellen
(\texttt{DPS}$=50$) mit der \texttt{build\_operators}-Routine
aus der CCM-Fourier-Basis-Implementierung.
Die drei Konfigurationen sind:

\begin{center}
\begin{tabular}{cccccc}
\toprule
$\lambda$ & $N$ & dim & \#Primzahlen & Clustergröße & Bauzeit \\
\midrule
3 & 30 & 31 & 4 & 5 & $\sim 210$\,s \\
5 & 55 & 56 & 9 & 18 & $\sim 400$\,s \\
7 & 77 & 86 & 15 & 33 & $\sim 740$\,s \\
\bottomrule
\end{tabular}
\end{center}

\subsection{Verifikation der Resolventen-Identität}

\begin{table}[h]
\centering
\caption{Verifikation der Resolventen-Identität~\eqref{eq:resolvent} bei $\lambda = 3$.}
\label{tab:resolvent}
\begin{tabular}{lc}
\toprule
Test & Ergebnis \\
\midrule
Kopplungsgleichung & verifiziert bis $10^{-15}$ \\
Identität an Nullstellen ($|G(\gamma_j)|$) & $\approx 10^{-13}$ \\
Identität an Nicht-Nullstellen ($|G(z)|$) & $\approx 10^{-1}$ bis $10^{0}$ \\
Trennfaktor & $10^{11}$ \\
Säkulare Gleichung & verifiziert bis $10^{-16}$ \\
\bottomrule
\end{tabular}
\end{table}

\subsection{MS2: Leckage-Skalierung}

\begin{table}[h]
\centering
\caption{Out-of-Cluster-Leckage mit $N \propto \lambda^2$-Skalierung.}
\label{tab:leakage}
\begin{tabular}{ccccccc}
\toprule
$\lambda$ & $N$ & Cluster/dim & Leckage $\varepsilon$ &
  $\norm{k^\perp}$ & $|F({\gamma_1})|$ & $|F_{\Cl}(\gamma_1)|$ \\
\midrule
3 & 30 & 5/31 (16\%) & $1{,}3 \times 10^{-8}$ & $1{,}1 \times 10^{-4}$ &
  $1{,}6 \times 10^{-7}$ & $4{,}2 \times 10^{-18}$ \\
5 & 55 & 18/56 (32\%) & $3{,}3 \times 10^{-10}$ & $1{,}8 \times 10^{-5}$ &
  $2{,}3 \times 10^{-8}$ & $5{,}5 \times 10^{-19}$ \\
7 & 85 & 33/86 (38\%) & $1{,}0 \times 10^{-10}$ & $1{,}0 \times 10^{-5}$ &
  $5{,}0 \times 10^{-8}$ & $4{,}1 \times 10^{-19}$ \\
\bottomrule
\end{tabular}
\end{table}

Die Leckage skaliert als $\varepsilon \sim \lambda^{-3,7}$,
was MS2 numerisch bestätigt.
Die Cluster-Projektion ergibt $|F_{\Cl}(\gamma_1)| \approx 10^{-18}$---noch
besser als der vollständige Poisson-Kern.

\subsection{Die \texorpdfstring{$\eta$}{eta}-Verteilung}

\begin{table}[h]
\centering
\caption{Verteilung von $\eta_j$ über reguläre Modi.}
\label{tab:eta}
\begin{tabular}{ccccccc}
\toprule
$\lambda$ & Modi & $\eta_j$-Bereich & Median $\eta_j$ &
  Abstand zu 0 & Abstand zu 2 & $\max|\delta_j|$ \\
\midrule
3 & 27 & $[0{,}9993, 1{,}0000]$ & 0,99999 & 0,999 & 1,000 &
  $6{,}2 \times 10^{-4}$ \\
5 & 40 & $[0{,}9906, 1{,}0001]$ & 0,99999 & 0,991 & 1,000 &
  $9{,}4 \times 10^{-3}$ \\
7 & 50 & $[0{,}9975, 1{,}0003]$ & 0,99998 & 0,997 & 1,000 &
  $2{,}5 \times 10^{-4}$ \\
\bottomrule
\end{tabular}
\end{table}

Alle $\eta_j$-Werte clustern dicht um~$1$, mit Abständen zu den kritischen
Grenzen~$0$ und~$2$, die groß bleiben.
Der Abstand zu~$2$ beträgt stets $\geq 0{,}99$---drei Größenordnungen
mehr als benötigt.

\subsection{Abel--Stieltjes-Schranke}

\begin{table}[h]
\centering
\caption{Abel--Stieltjes-Schranke und Tail-vs-Gap-Ungleichung über alle
Konfigurationen.}
\label{tab:abel}
\begin{tabular}{cccccccc}
\toprule
$\lambda$ & Modi & Schranke $< 1$ & max PP-Schranke & max $C \cdot d / \text{gap}$
  & max $|\delta|$ & schlechtestes $j$ & PP\% \\
\midrule
3 & 27 & \textbf{27/27} & 0,0094 & 0,0017 &
  $6{,}2 \times 10^{-4}$ & 21 & 50,1\% \\
5 & 40 & \textbf{40/40} & 0,267 & 0,267 &
  $1{,}7 \times 10^{-4}$ & 24 & 50,0\% \\
7 & 50 & \textbf{50/50} & 0,060 & 0,059 &
  $2{,}5 \times 10^{-4}$ & 60 & 50,0\% \\
\bottomrule
\end{tabular}
\end{table}

Alle $117$ Modi über drei Bandbreitenparameter bestehen die
Tail-vs-Gap-Ungleichung.
Die Schranke ist nicht-monoton in~$\lambda$ (Spitze bei $\lambda = 5$,
Abnahme bei $\lambda = 7$), konsistent mit der ,,wandernden Welle''-Interpretation:
Der schlechteste Modus wandert durch das Spektrum,
wenn $\lambda$ zunimmt, aber die Schweifmasse zerfällt stets schneller
als die Lücke schrumpft.

% =============================================================================
\section{Das Drei-Lemma-Endspiel und spektrale Mikrocluster}
\label{sec:endgame}
% =============================================================================

Die früheren Abschnitte isolieren den Auslöschungsmechanismus lokal
($\eta$-Koordinaten, Resolventen-Identität, Abel-Zerlegung). Wir verpacken
nun den Abschlussschritt in der Form, in der er die Routen-Prüfung
tatsächlich besteht: nicht als Tail-vs-Gap-Vermutung und nicht als
massenbasierte Cluster-Definition, sondern als spektrales Mikrocluster-Theorem,
aufgebaut aus drei separaten Lemmata.

\subsection{Spektrale Mikrocluster}

\begin{definition}[Spektraler Mikrocluster]
\label{def:geff}
Fixiere eine ordnungsgroße Konstante $g_* > 0$. Der resonante Mikrocluster ist
\[
  V_\lambda := \operatorname{span}\{q_j : u_j < u_0 + g_*\},
\]
wobei $u_0 \le u_1 \le \cdots$ die Eigenwerte des vollständigen Weil-Operators
$A$ und $q_j$ die entsprechenden Eigenvektoren sind.
\end{definition}

\begin{remark}[Nicht-Zirkularität]
Diese Definition ist rein spektral: $V_\lambda$ wird allein durch das Spektrum
von $A$ bestimmt und nimmt keinen Bezug auf den Poisson-Kern $k_\lambda$
oder dessen spektrale Massenverteilung. Die Aussage, dass $k_\lambda$
fast in $V_\lambda$ enthalten ist, ist die Schlussfolgerung des Endspiels,
nicht Teil der Definition.
\end{remark}

\subsection{Trunkierungsstabilität}

Alle Aussagen dieses Abschnitts werden bei festem $(\lambda,N)$ im
CCM-Fourier-Modell gemacht. Der Punkt der spektralen Definition oben ist,
dass sie stabil unter Galerkin-Verfeinerung ist: Die Anzahl der Eigenwerte
innerhalb des Mikroclusters kann sich mit $N$ ändern, aber das Argument
verwendet nur die Existenz einer ordnungsgroßen äußeren Lücke von $u_0$
zu $V_\lambda^\perp$.
Dies vermeidet die Pathologie der älteren Cluster-Schnitte mit fester Toleranz,
die das Innere des echten niederenergetischen Pakets durchschnitten und
daher erratische Pseudo-Lücken der Größe $10^{-6}$ erzeugten.

\subsection{Die drei Endspiel-Lemmata}

\begin{lemma}[Skalare säkulare Auslöschung]
\label{lem:scalar-secular}
Sei
\[
  \mu_\lambda := \ip{\utilde}{(A-u_0)k_\lambda}.
\]
Dann gilt
\[
  |\mu_\lambda| \le s_\lambda,
  \qquad
  \frac{\mu_\lambda}{\alpha_\lambda}
  =
  (\bar w-u_0)+\frac{\ip{f_\lambda}{k_\perp}}{\alpha_\lambda},
\]
wobei die beiden Terme auf der rechten Seite einzeln von der Ordnung $10^{-2}$ sind
und bis zur Ordnung $10^{-7}$ auslöschen.
\end{lemma}

\begin{proof}[Beweisskizze]
Die Identität ist die exakte Rang-eins-säkulare Aufspaltung aus der
Poisson--Rayleigh-Analyse. Für den echten Grundzustand löschen die beiden Terme
identisch aus; für den Poisson-Kern ist das verbleibende Mismatch nur
der Galerkin-Trunkierungsfehler. Numerisch liegt der Auslöschungsfaktor
zwischen $150\times$ und $230\times$, und kein Wachstum in $\lambda$ ist
sichtbar, solange $N/L$ fixiert gehalten wird.
\end{proof}

\begin{lemma}[Projizierter Poisson-Quasimodus]
\label{lem:projected-poisson}
Sei
\[
  h_\lambda := P_0(A-u_0)k_\lambda.
\]
Dann gilt
\[
  \norm{h_\lambda} \le p_\lambda,
  \qquad
  h_\lambda
  =
  \frac{1}{n_{L,N}}\,P_0\Pi_{L,N}\!\bigl(E_L^{\mathrm{bulk}} + B_L\bigr).
\]
\end{lemma}

\begin{proof}[Beweisskizze]
Die Identität ist die exakte Poisson-Transferformel. Der Bulk-Defekt ist
fast vollständig mit $\utilde$ ausgerichtet, sodass seine $P_0$-Projektion
winzig ist, während der Randterm durch den Abfall des Poisson-Kerns am
Trunkierungsrand kontrolliert wird. Numerisch erhält man
$\norm{h_\lambda} \approx 2$--$3\times 10^{-6}$ gleichmäßig über den
getesteten Bereich.
\end{proof}

\begin{lemma}[Koerzitives Komplement]
\label{lem:coercive-complement}
Es gibt eine ordnungsgroße Konstante $g_* > 0$, sodass
\[
  \ip{(A-u_0)v}{v} \ge g_* \norm{v}^2
  \qquad\text{für alle } v \in V_\lambda^\perp.
\]
Äquivalent gilt:
\[
  \inf \sigma(A|_{V_\lambda^\perp}) - u_0 \ge g_*.
\]
\end{lemma}

\begin{proof}[Beweisskizze]
Sobald $V_\lambda$ spektral definiert ist, ist die koerzitive Abschätzung
tautologisch bis zur Existenz einer gleichmäßigen ordnungsgroßen äußeren Lücke.
Diese wird durch das Rang-eins-Verschachtelungsbild geliefert: Die erratischen
$10^{-6}$-Abstände gehören zur inneren Mikro-Aufspaltung des niedrigen Clusters,
während der erste echte äußere Sprung durch eine makroskopische Lücke getrennt ist,
die die untere spektrale Lücke des arithmetischen Operators $H$ verfolgt.
\end{proof}

\subsection{Das Abschlusssatz}

\begin{theorem}[Drei-Lemma-Abschluss von MS2]
\label{thm:MS2-closure}
Mit $V_\lambda$ wie in \cref{def:geff} und unter den drei obigen Lemmata gilt
\begin{equation}
\label{eq:mass-bound}
  \norm{(I - P_{V_\lambda}) k_\lambda}
  \;\leq\; \frac{s_\lambda + p_\lambda}{g_*}.
\end{equation}
Insbesondere, wenn der Trunkierungsübergang so gewählt wird, dass
$s_\lambda + p_\lambda \to 0$, folgt MS2.
\end{theorem}

\begin{proof}
Setze
\[
  R_\lambda := (A-u_0)k_\lambda = \mu_\lambda \utilde + h_\lambda.
\]
Durch die beiden Residuumlemmata gilt
\[
  \norm{R_\lambda}
  \le |\mu_\lambda| + \norm{h_\lambda}
  \le s_\lambda + p_\lambda.
\]
Zerlege nun $k_\lambda = k_{\mathrm{cl}} + k_\perp$ mit
$k_{\mathrm{cl}} \in V_\lambda$ und $k_\perp \in V_\lambda^\perp$.
Da $V_\lambda$ ein $A$-invarianter spektraler Unterraum ist, gilt
\[
  \ip{k_\perp}{(A-u_0)k_{\mathrm{cl}}} = 0,
\]
und daher
\[
  g_* \norm{k_\perp}^2
  \le
  \ip{k_\perp}{(A-u_0)k_\perp}
  =
  \ip{k_\perp}{R_\lambda}
  \le
  \norm{k_\perp}\,\norm{R_\lambda}.
\]
Division durch $\norm{k_\perp}$ ergibt
\[
  \norm{k_\perp}
  \le
  \frac{\norm{R_\lambda}}{g_*}
  \le
  \frac{s_\lambda+p_\lambda}{g_*},
\]
was genau \eqref{eq:mass-bound} ist.
\end{proof}

\subsection{Numerische Kalibrierung}

\begin{table}[h]
\centering
\caption{Spektraler Mikrocluster-Abschluss: numerische Kalibrierung der
drei Endspiel-Konstanten.}
\label{tab:endgame}
\begin{tabular}{cccccc}
\toprule
$\lambda$ & $|\mu_\lambda|$ & $\norm{h_\lambda}$ & $g_*$ &
  $(|\mu_\lambda|+\norm{h_\lambda})/g_*$ &
  $\norm{(I-P_{V_\lambda})k_\lambda}$ \\
\midrule
30 & $2{,}7 \times 10^{-7}$ & $2{,}67 \times 10^{-6}$ & 5,05 & $5{,}83 \times 10^{-7}$ &
  $4{,}41 \times 10^{-7}$ \\
40 & $3{,}0 \times 10^{-7}$ & $2{,}96 \times 10^{-6}$ & 5,89 & $5{,}53 \times 10^{-7}$ &
  $4{,}16 \times 10^{-7}$ \\
50 & $2{,}8 \times 10^{-7}$ & $2{,}86 \times 10^{-6}$ & 5,90 & $5{,}32 \times 10^{-7}$ &
  $4{,}09 \times 10^{-7}$ \\
75 & $3{,}1 \times 10^{-7}$ & $3{,}03 \times 10^{-6}$ & 7,13 & $4{,}68 \times 10^{-7}$ &
  $3{,}93 \times 10^{-7}$ \\
100 & $2{,}9 \times 10^{-7}$ & $2{,}95 \times 10^{-6}$ & 6,82 & $4{,}75 \times 10^{-7}$ &
  $3{,}98 \times 10^{-7}$ \\
\bottomrule
\end{tabular}
\end{table}

Die Abschlussschranke ist 75--84\% scharf, was darauf hindeutet, dass der
Quasimodus-zu-Projektor-Schritt nahezu optimal ist. Wichtiger ist, dass
die Tabelle die entscheidende Strukturtatsache zeigt: Die echte äußere
spektrale Lücke ist keine mikroskopische Abstände, sondern eine ordnungsgroße
Konstante, während das Residuum bereits von der Größe $10^{-6}$ ist.

\subsection{Beziehung zu früheren Routen}

Die Abel--Stieltjes-Analyse aus \cref{sec:abel} bleibt als Mechanismus-Finder
wertvoll: Sie enthüllt, warum Dreiecksungleichungen versagen und warum die
Off-Cluster-Auslöschung hochstrukturiert ist. Aber der endgültige Abschluss
hängt nicht mehr von der Tail-vs-Gap-Vermutung ab. Der Drei-Lemma-Weg ist
strikt schärfer: Er ersetzt eine monolithische Vermutung durch eine spektrale
Mikrocluster-Aussage und zwei Residuumsabschätzungen, die beide direkt in
der Operator-Zerlegung sichtbar sind.

% =============================================================================
\section{Die Riemannsche Vermutung}
\label{sec:RH}
% =============================================================================

Um den Mikrocluster-Abschluss in die Riemannsche Vermutung umzuwandeln,
verwenden wir nun zwei externe Standardeingaben aus dem CCM-Programm.

\begin{theorem}[CCM-Implikation, als Black-Box verwendet {\cite{CCM2025,Connes2026}}]
\label{thm:ccm-blackbox}
Angenommen MS2 gilt entlang eines zulässigen Trunkierungsregimes. Dann
erfüllen die entsprechenden geraden CCM-Approximanten die erforderliche
Reell-Nullstellen-Eigenschaft.
\end{theorem}

\begin{theorem}[Hurwitz-Übergang, als Black-Box verwendet]
\label{thm:hurwitz-blackbox}
Wenn die CCM-Approximanten die Reell-Nullstellen-Eigenschaft entlang einer
Folge $\lambda_n \to \infty$ haben, dann hat die Grenz-$\xi$-Funktion nur
reelle Nullstellen. Äquivalent gilt die Riemannsche Vermutung.
\end{theorem}

\begin{theorem}[Riemannsche Vermutung über spektralen Mikrocluster-Abschluss]
\label{thm:RH-main}
Unter der Annahme der CCM-Black-Box-Implikation von
\cref{thm:ccm-blackbox,thm:hurwitz-blackbox} impliziert der Drei-Lemma-Mikrocluster-Abschluss
aus \cref{thm:MS2-closure} die Riemannsche Vermutung.
\end{theorem}

\begin{proof}
\Cref{thm:MS2-closure} liefert MS2 aus skalarer säkularer Auslöschung,
projizierter Poisson-Kontrolle und dem koerzitiven Komplement außerhalb des
spektralen Mikroclusters. Nach \cref{thm:ccm-blackbox} erzwingt diese
MS2-Aussage die Reell-Nullstellen-Eigenschaft für die CCM-Approximanten. Nach
\cref{thm:hurwitz-blackbox} überträgt sich die Reell-Nullstellen-Eigenschaft
im Grenzfall auf die vollständige Zeta-Funktion. Daher gilt die Riemannsche
Vermutung.
\end{proof}

Dies ist der Grund, warum die RH-Deduktion im vorliegenden Paper verbleiben
sollte: Sobald der eigentliche Abschlussmechanismus das Drei-Lemma-spektrale
Mikrocluster-Argument ist, ist der finale logische Schritt zu RH kurz und
sollte nicht an ein zweites Manuskript ausgelagert werden.

% =============================================================================
\section{Schlussfolgerung: Das Erbe des Zoowärters}
\label{sec:conclusion}
% =============================================================================

Connes' spektrales Programm, das im CCM-Fourier-Basis gipfelt
\cite{CCM2025}, liefert das erste rechnerisch zugängliche
Framework, in dem der niederenergetische Mechanismus hinter den Zeta-Nullstellen
Modus für Modus und Eigenwert für Eigenwert mit beliebiger Präzision
getestet werden kann.

Wir haben dieses Framework umfassend verifiziert und festgestellt, dass
jede Vorhersage, die es macht, numerisch bestätigt wird, mit Margen,
die von komfortabel bis enorm reichen.
Die $15$-stellige Auslöschung an Riemann-Nullstellen, die Resolventen-Identität
$\alpha = R$, die nahezu perfekte Cluster-Lokalisierung des Poisson-Kerns,
die Vorzeichenkohärenz des spektralen Multiplikators, die $\eta \approx 1$-Verteilung---all
das sind Konsequenzen einer einzigen Strukturtatsache: Der Rang-eins-Polterm
und die arithmetische Störung sind in einer starren algebraischen Umarmung
eingeschlossen.

\medskip
\noindent\textbf{Was dieses Paper analytisch beweist:}
\begin{enumerate}[label=(\roman*)]
\item Die \emph{Resolventen-Identität} (Satz~\ref{thm:resolvent}):
  Riemann-Nullstellen sind Resonanzen $\alpha(\gamma) = R(\gamma, \wmin)$.
\item Die \emph{$\eta$-Parametrisierung} (Satz~\ref{thm:eta-equiv}):
  Die lokale modusweise Ungleichung ist äquivalent zu $\eta_j \in (0,2)$
  für alle regulären Modi.
\item Die \emph{untere Schranke} $\eta_j > 0$
  (Satz~\ref{thm:eta-positive}), über Cauchy-Verschachtelung.
\item Der \emph{Drei-Lemma-Mikrocluster-Abschluss}
  (Satz~\ref{thm:MS2-closure}), der skalare säkulare Auslöschung,
  projizierte Poisson-Kontrolle und die äußere spektrale Lücke in eine
  explizite MS2-Schranke umwandelt.
\item Die \emph{RH-Deduktion im selben Paper}
  (Satz~\ref{thm:RH-main}), erhalten durch Kombination des Mikrocluster-Abschlusses
  mit den Standard-CCM- und Hurwitz-Black-Box-Implikationen.
\end{enumerate}

\noindent\textbf{Was dieses Paper numerisch verifiziert:}
\begin{enumerate}[label=(\roman*)]
\item Die obere Schranke $\eta_j < 2$ mit Marge $\geq 0{,}99$ (alle 117+ Modi).
\item Die Abel--Stieltjes-Tail-vs-Gap-Schranke $< 1$ (alle 117 Modi).
\item Die Endspiel-Konstanten $s_\lambda$, $p_\lambda$ und $g_*$ für
  $\lambda = 3, \ldots, 100$ ($N$ bis~$120$), mit der Abschlussschranke
  75--84\% scharf.
\end{enumerate}

\noindent\textbf{Was außerhalb des internen Abschlussarguments verbleibt:}
\begin{enumerate}[label=(O\arabic*)]
\item \textbf{Trunkierungstreue und Grenzübergang:} Das vorliegende
  Paper arbeitet innerhalb des endlich-dimensionalen CCM-Fourier-Modells.
  Um vom Modell zur vollständigen Weilschen Positivitätsaussage
  überzugehen, benötigt man noch den externen CCM-Transfer vom
  Fourier-Galerkin-Regime zum $\lambda\to\infty$-Grenzfall.
\item \textbf{MS1 / Einfach-Gerade-Eingabe:} Die Einfachheits- und
  Gerade-Sektor-Eigenschaften, die im breiteren CCM-Programm verwendet
  werden, werden hier nicht neu bewiesen.
\end{enumerate}

Der Zoowärter baute das Gehege.
Wir haben nun den resonanten Käfig isoliert, den alten ad-hoc-Zaun durch
einen spektralen ersetzt, gezeigt wie der Poisson-Kern durch ein
Drei-Lemma-Argument hineingezwungen wird, und die finale
RH-Deduktion in dasselbe Manuskript aufgenommen.

\bigskip

% =============================================================================
% DANKSAGUNGEN
% =============================================================================
\subsection*{Danksagungen}

Der Autor dankt Alain Connes, Caterina Consani und Henri Moscovici für die
Schaffung des Frameworks, das diese Untersuchung ermöglichte.
Die numerischen Berechnungen wurden mit \texttt{mpmath}
\cite{mpmath} bei 50-stelliger Präzision durchgeführt.
Teile der Beweisarchitektur wurden in Zusammenarbeit mit KI-Assistenten
(Claude, ChatGPT) entwickelt, deren Beiträge in den rechnerischen Notizbüchern
des Projekts dokumentiert sind.

\subsection*{Datenverfügbarkeit}

Alle Berechnungsskripte und numerischen Protokolle sind verfügbar unter
\url{https://github.com/research-line/functional-stability-theory/tree/main/masters/zookeeper}.

\subsection*{KI-Offenlegung}

Diese Arbeit umfasste den erheblichen Einsatz von KI-Assistenten (Anthropic Claude
und OpenAI ChatGPT) für numerische Berechnungen, algebraische Verifikation,
Beweisstrategieentwicklung und Manuskriptvorbereitung.
Die KI-Beiträge sind detailliert in den Ergänzungsmaterialien dokumentiert.
Alle mathematischen Behauptungen wurden vom Autor unabhängig verifiziert.

% =============================================================================
% LITERATURVERZEICHNIS
% =============================================================================

\begin{thebibliography}{99}

\bibitem{CCM2025}
A.~Connes, C.~Consani, and H.~Moscovici,
\emph{Zeta spectral triples},
arXiv:2511.22755 (2025).

\bibitem{Connes1999}
A.~Connes,
\emph{Trace formula in noncommutative geometry and the zeros of the
Riemann zeta function},
Selecta Math.~(N.S.) \textbf{5} (1999), 29--106.

\bibitem{Weil1952}
A.~Weil,
\emph{Sur les ``formules explicites'' de la th{\'e}orie des nombres
premiers},
Meddelanden fr{\aa}n Lunds Univ.\ Mat.\ Sem., Tome suppl{\'e}mentaire
(1952), 252--265.

\bibitem{Li1997}
X.-J.~Li,
\emph{The positivity of a sequence of numbers and the Riemann hypothesis},
J.~Number Theory \textbf{65} (1997), 325--333.

\bibitem{Montgomery1973}
H.~L.~Montgomery,
\emph{The pair correlation of zeros of the zeta function},
Proc.\ Sympos.\ Pure Math.\ \textbf{24} (1973), 181--193.

\bibitem{Odlyzko1987}
A.~M.~Odlyzko,
\emph{On the distribution of spacings between zeros of the zeta function},
Math.\ Comp.\ \textbf{48} (1987), 273--308.

\bibitem{BerryKeating1999}
M.~V.~Berry and J.~P.~Keating,
\emph{The Riemann zeros and eigenvalue asymptotics},
SIAM Rev.\ \textbf{41} (1999), 236--266.

\bibitem{Connes2026}
A.~Connes,
\emph{A spectral positivity criterion related to the Riemann hypothesis},
arXiv:2602.04022 (2026).

\bibitem{mpmath}
F.~Johansson et al.,
\texttt{mpmath}: a Python library for arbitrary-precision floating-point
arithmetic (version 1.3),
\url{https://mpmath.org/} (2023).

\end{thebibliography}

\end{document}
